% problem_id: S001-lemma-affineness-of-large-open
% source_id: S001 (Stacks Project)
% source_file: more-morphisms.tex
% statement_locator: more-morphisms.tex lines 22965--22975
% proof_locator: more-morphisms.tex lines 22977--23077
% env: lemma
% id_note: label 在本来源内唯一
% pairing: tight (verified)
% status: complete (statement + author proof verbatim + reason + locators)
%
\begin{lemma}
\label{lemma-affineness-of-large-open}
Let $f : X \to Y$ be a morphism of affine schemes which is
quasi-finite and of finite presentation.
Let $w : X \to \mathbf{Z}_{> 0}$ be a positive weighting of $f$.
Let $d < \infty$ be the maximum value of $\int_f w$. The open
$$
Y_d = \{y \in Y \mid (\textstyle{\int}_f w)(y) = d \}
$$
of $Y$ is affine.
\end{lemma}

\begin{proof}
Observe that $\int_f w$ attains its maximum by Lemma \ref{lemma-max-int-w}.
The set $Y_d$ is open by Lemma \ref{lemma-semicontinuous-int-w}.
Thus the statement of the lemma makes sense.

\medskip\noindent
Reduction to the Noetherian case; please skip this paragraph.
Recall that a finite type morphism is quasi-finite if and only
if it has relative dimension $0$, see
Morphisms, Lemma \ref{morphisms-lemma-locally-quasi-finite-rel-dimension-0}.
By Lemma \ref{lemma-Noetherian-approximation-dimension-d}
applied with $d = 0$ we can find a quasi-finite morphism $f_0 : X_0 \to Y_0$
of affine Noetherian schemes and a morphism $Y \to Y_0$ such that $f$
is the base change of $f_0$. Then we can write $Y = \lim Y_i$ as a directed
limit of affine schemes of finite type over $Y_0$, see
Algebra, Lemma \ref{algebra-lemma-ring-colimit-fp}.
By Lemma \ref{lemma-descend-weighting}
we can find an $i$ such that our weighting $w$
descends to a weighting $w_i$ of the base change $f_i : X_i \to Y_i$
of $f_0$. Now if the lemma holds for $f_i, w_i$, then it implies
the lemma for $f$ as formation of $\int_f w$ commutes with base
change, see Lemma \ref{lemma-weighting-check-after-etale-base-change}.

\medskip\noindent
Assume $X$ and $Y$ Noetherian. Let $X' \to Y'$ be the base change of $f$
by a morphism $g : Y' \to Y$. The formation of $\int_f w$ and hence
the open $Y_d$ commute with base change. If $g$ is finite and surjective, then
$Y'_d \to Y_d$ is finite and surjective. In this case proving that
$Y_d$ is affine is equivalent to showing that $Y'_d$ is affine, see
Cohomology of Schemes, Lemma
\ref{coherent-lemma-image-affine-finite-morphism-affine-Noetherian}.

\medskip\noindent
We may choose an immersion $X \to T$ with $T$ finite over $Y$, see Lemma
\ref{lemma-quasi-finite-separated-pass-through-finite}.
We are going to apply Morphisms, Lemma \ref{morphisms-lemma-massage-finite}
to the finite morphism $T \to Y$. This lemma tells us that
there is a finite surjective morphism $Y' \to Y$ such that
$Y' \times_Y T$ is a closed subscheme of a scheme $T'$ finite over $Y'$
which has a special form.
By the discussion in the first paragraph, we may replace $Y$ by $Y'$,
$T$ by $T'$, and $X$ by $Y' \times_Y X$.
Thus we may assume there is an immersion $X \to T$ (not necessarily
open or closed) and closed subschemes
$T_i \subset T$, $i = 1, \ldots, n$ where
\begin{enumerate}
\item $T \to Y$ is finite (and locally free),
\item $T_i \to Y$ is an isomorphism, and
\item $T = \bigcup_{i = 1, \ldots, n} T_i$ set theoretically.
\end{enumerate}
Let $Y' = \coprod Y_k$ be the disjoint union of the irreducible
components of $Y$ (viewed as integral closed subschemes of $Y$).
Then we may base change once more by $Y' \to Y$; here we
are using that $Y$ is Noetherian. Thus we may in
addition assume $Y$ is integral and Noetherian.

\medskip\noindent
We also may and do assume that $T_i \not = T_j$ if $i \not = j$ by
removing repeats. Since $Y$ and hence all $T_i$ are integral, this
means that if $T_i$ and $T_j$ intersect, then they intersect in a
closed subset which maps to a proper closed subset of $Y$.

\medskip\noindent
Observe that $V_i = X \cap T_i$ is a locally closed subset
which is in addition a closed subscheme of $X$ hence affine.
Let $\eta \in Y$ and $\eta_i \in T_i$ be the generic points.
If $\eta \not \in Y_d$, then $Y_d = \emptyset$ and we're done.
Assume $\eta \in Y_d$. Denote $I \in \{1, \ldots, n\}$
the subset of indices $i$ such that $\eta_i \in V_i$.
For $i \in I$ the locally closed subset $V_i \subset T_i$
contains the generic point of the irreducible space $T_i$
and hence is open. On the other hand, since $f$ is open
(Lemma \ref{lemma-weighting-universally-open}),
for any $x \in X$ we can find an $i \in I$ and a specialization
$\eta_i \leadsto x$. It follows that $x \in T_i$ and hence
$x \in V_i$. In other words, we see that $X = \bigcup_{i \in I} V_i$
set theoretically.
We claim that $Y_d = \bigcap_{i \in I} \Im(V_i \to Y)$; this will
finish the proof as the intersection of affine opens
$\Im(V_i \to Y)$ of $Y$ is affine.

\medskip\noindent
For $y \in Y$ let $f^{-1}(\{y\}) = \{x_1, \ldots, x_r\}$ in $X$.
For each $i \in I$ there is at most one $j(i) \in \{1, \ldots, x_r\}$
such that $\eta_i \leadsto x_{j(i)}$. In fact, $j(i)$ exists and is
equal to $j$ if and only if $x_j \in V_i$. If $i \in I$ is such that
$j = j(i)$ exists, then $V_i \to Y$ is an isomorphism in a neighbourhood
of $x_j \mapsto y$. Hence $\bigcup_{i \in I,\ j(i) = j} V_i \to Y$
is finite after replacing source and target by neighbourhoods of
$x_j \mapsto y$. Thus the definition of a weighting tells us that
$w(x_j) = \sum_{i \in I,\ j(i) = j} w(\eta_i)$.
Thus we see that
$$
(\textstyle{\int}_f w)(\eta) =
\sum\nolimits_{i \in I} w(\eta_i) \geq
\sum\nolimits_{j(i)\text{ exists}} w(\eta_i) =
\sum\nolimits_j w(x_j) = (\textstyle{\int}_f w)(y)
$$
Thus equality holds if and only if $y$ is contained in
$\bigcap_{i \in I} \Im(V_i \to Y)$ which is what we wanted to show.
\end{proof}
